%% =================================================================
%% Wei, Mei, Zhao, Xiao, Sun, Zhang, Guo.
%% "Nondestructively identifying the mechanical behavior of soft
%%  tissues using surface deformation with an explicit inverse
%%  approach."
%% Appl. Math. Modelling 134 (2024) 126--147.
%%
%% Reconstruction of page 9 (printed p. 134) as study notes.
%% This page is almost entirely figures (Fig. 7 and Fig. 8), with
%% minimal connecting prose at the top and bottom.
%% Prose: extracted verbatim from the PDF text layer via pdftotext.
%% Math:  none on this page.
%% Figures 7 & 8: placeholder boxes with captions + visual descriptions.
%% =================================================================

\documentclass[11pt]{article}

\usepackage[margin=1in]{geometry}
\usepackage{amsmath, amssymb}
\usepackage{mathrsfs}
\usepackage{bm}
\usepackage{fancyhdr}

\pagestyle{fancy}
\fancyhf{}
\lhead{\small Y.~Mei et al.}
\rhead{\small Applied Mathematical Modelling 134 (2024) 126--147}
\cfoot{\small \thepage}
\renewcommand{\headrulewidth}{0.4pt}

\setcounter{page}{134}

\begin{document}

% ---------- Figure 7 placeholder ----------
\noindent
\begin{center}
\fbox{\parbox{0.92\textwidth}{\centering\vspace{1em}
\textbf{[Figure 7 — placeholder]}\\[0.8em]
Six-panel grid showing the optimization trajectory for the
explicit inverse scheme on \textbf{Sample 1} (the homogeneous case,
target $\text{SM} = 3.00$\,MPa), starting from a 6-layer initial
guess. Each panel is a rectangular shear-modulus distribution
($1.0$\,mm $\times 0.2$\,mm) with a jet colorbar.\\[0.4em]
\begin{tabular}{@{}ll@{}}
\textbf{(a) Initial guess} & uniform blue, SM $= 1.00$\,MPa everywhere \\
\textbf{(b) Iteration 5} & SM range 1.27--1.95, layered artifacts emerging \\
\textbf{(c) Iteration 12} & SM range 1.50--3.11, layered structure developed \\
\textbf{(d) Iteration 24} & SM range 1.63--3.03, approaching target \\
\textbf{(e) Iteration 48} & SM range 2.97--3.00, nearly converged \\
\textbf{(f) Iteration 52 (Final)} & uniform red, SM $= 3.00$\,MPa everywhere
\end{tabular}\\[0.4em]
The trajectory shows the 6-layer parameterisation collapsing to a
uniform solution over 52 iterations --- because the target is
homogeneous, the optimiser correctly drives all six layers to the
same shear modulus.
\vspace{1em}}}
\end{center}
\vspace{0.3em}
\begin{center}
\textbf{Fig. 7.} The optimization process with six layers initially
utilizing the explicit inverse scheme for Sample 1 (Without noise,
Unit: MPa).
\end{center}

\vspace{0.8em}

% ---------- Figure 8 placeholder ----------
\noindent
\begin{center}
\fbox{\parbox{0.92\textwidth}{\centering\vspace{1em}
\textbf{[Figure 8 — placeholder]}\\[0.8em]
Six-panel side-by-side comparison of the
\textbf{Explicit Inverse Method} (left column) vs.\
\textbf{Implicit Inverse Method} (right column) on Sample 1
(homogeneous, target SM $= 3.00$\,MPa), with 3, 6, and 9 surface
displacement datasets as inputs.\\[0.4em]
\textbf{Explicit (left column)}: all three rows show a perfectly
uniform shear-modulus field at SM $= 3.00$\,MPa (uniform blue,
colorbar pinned 3.00--3.00). The explicit parameterisation
recovers the homogeneous target exactly.\\[0.4em]
\textbf{Implicit (right column)}: all three rows show tiny spatial
artifacts around the true value of 3.00\,MPa --- (d) 3 fields: SM
range 2.9995--3.0001; (e) 6 fields: SM range 2.9992--3.0001;
(f) 9 fields: SM range 3.0000--3.0002. The color blobs are
visible even though the numerical error is on the order of
$10^{-4}$--$10^{-3}$\,MPa.\\[0.4em]
\textbf{Takeaway}: the explicit method produces a perfectly uniform
field on the homogeneous case; the implicit method is quantitatively
accurate but shows residual spatial noise patterns that decrease
(slightly) as more displacement datasets are incorporated.
\vspace{1em}}}
\end{center}
\vspace{0.3em}
\begin{center}
\textbf{Fig. 8.} The reconstructed results with varying 3, 6 and 9
surface displacement datasets corresponding to Sample 1 (Without
noise, Unit: MPa).
\end{center}

\vspace{1em}

% ---------- Continuation prose (from page 8) + noise-test paragraph ----------
\noindent
inverse method shows a slight spatial variation, which can be
enhanced by incorporating more surface datasets.

After adding 1\,\% noise to the surface displacement datasets, we
observed that the shear modulus distributions were accurately mapped
for both inverse methods in the homogeneous case (refer to Fig.~11).
The average relative error for both implicit and explicit inverse
methods was less than 0.2\,\% (see Table~1). However, when it came
to Sample 2 and 3 (see Fig.~12, Fig.~13), the interface between
neighboring layers in the implicit inverse method's results was not
as clear as the target. As the total number of layers increases, the
challenges in solving the inverse problems escalate accordingly.
Consequently, the relative error for Sample 2 and Sample 3 exhibits
a lower quality compared to Sample 1 (see Table~1). Besides, the
average relative error for the implicit inverse method was

% [page ends here — continues on page 10]

\end{document}
